% Section 15.4: physical PDF pp. 37--41 (printed pp. 14--18).
% Semantic coverage: begins with the heading on p. 37 and ends immediately
% before Section 15.5. Inventory: Eqs. (15.4.1)--(15.4.17), one footnote,
% one citation marker [9], and one centered * * * divider. Uncertainty: none.

\subsection{Quantization}

We now proceed to quantize the gauge theories described in the previous two
sections. The Lagrangian density is taken in the form \eqref{eq:15.3.1}:
\begin{equation}
  \mathcal L=-\frac14 F_{\alpha\mu\nu}F_\alpha{}^{\mu\nu}
  +\mathcal L_M(\psi,D_\mu\psi),
  \tag{15.4.1}\label{eq:15.4.1}
\end{equation}
with
\begin{align*}
  F_{\alpha\mu\nu}&\equiv\partial_\mu A_{\alpha\nu}
  -\partial_\nu A_{\alpha\mu}+C_{\alpha\beta\gamma}
  A_\beta{}_{\mu}A_\gamma{}_{\nu},\\
  D_\mu\psi&\equiv\partial_\mu\psi-it_\alpha A_{\alpha\mu}\psi.
\end{align*}
We cannot immediately quantize this theory by setting commutators equal to
$i$ times the corresponding Poisson brackets. The problem is one of
constraints. In the terminology of Dirac, described in Section 7.6, there is
a primary constraint that
\begin{equation}
  \Pi_{\alpha0}\equiv\frac{\partial\mathcal L}
  {\partial(\partial_0A_\alpha{}^0)}=0
  \tag{15.4.2}\label{eq:15.4.2}
\end{equation}
and a secondary constraint provided by the field equation for $A_\alpha{}^0$:
\begin{align}
  -\partial_\mu\frac{\partial\mathcal L}
  {\partial(\partial_\mu A_{\alpha0})}
  +\frac{\partial\mathcal L}{\partial A_{\alpha0}}
  &=\partial_\mu F_\alpha{}^{\mu0}
  +F_\gamma{}^{\mu0}C_{\gamma\alpha\beta}A_{\beta\mu}+J_\alpha{}^0
  \notag\\
  &=\partial_k\Pi_\alpha{}^k
  +\Pi_\gamma{}^k C_{\gamma\alpha\beta}A_{\beta k}+J_\alpha{}^0=0,
  \tag{15.4.3}\label{eq:15.4.3}
\end{align}
where $\Pi_\alpha{}^k\equiv\partial\mathcal L/
\partial(\partial_0A_{\alpha k})=F_\alpha{}^{k0}$ is the `momentum'
conjugate to $A_{\alpha k}$, with $k$ running over the values 1, 2, 3. The
Poisson brackets of $\Pi_{\alpha0}$ and
$\partial_k\Pi_\alpha{}^k+\Pi_\gamma{}^kC_{\gamma\alpha\beta}A_{\beta k}
+J_\alpha{}^0$ vanish (because the latter quantity is independent of
$A_\alpha{}^0$), so these are first class constraints, which cannot be dealt
with by replacing Poisson brackets with Dirac brackets.

As in the case of electrodynamics, we deal with these constraints by choosing
a gauge. The Coulomb gauge adopted for electrodynamics would lead to painful
complications here,\footnote{In addition to purely algebraic complications,
Coulomb gauge (like many other gauges) has a problem known as the Gribov
ambiguity~\hyperref[ch15-ref-9]{[9]}: even with the condition that $A_\alpha$
vanishes at spatial infinity, for each solution of the Coulomb gauge condition
$\boldsymbol\nabla\cdot\mathbf A_\alpha=0$ there are other solutions that differ
by finite gauge transformations. The Gribov ambiguity will not bother us here,
because we quantize in axial gauge where it is absent, and we shall use other
gauges like Lorentz gauge only to generate a perturbation series.} so instead
we will work in what is known as \emph{axial gauge}, based on the condition
\begin{equation}
  A_{\alpha3}=0.
  \tag{15.4.4}\label{eq:15.4.4}
\end{equation}
The canonical variables of the gauge field are then $A_{\alpha i}$, with $i$
now running over the values 1 and 2, together with their canonical conjugates
\begin{equation}
  \Pi_{\alpha i}\equiv\frac{\partial\mathcal L}
  {\partial(\partial_0A_{\alpha i})}=-F_\alpha{}^{0i}
  =\partial_0A_{\alpha i}-\partial_iA_{\alpha0}
  +C_{\alpha\beta\gamma}A_\beta{}^0A_{\gamma i}.
  \tag{15.4.5}\label{eq:15.4.5}
\end{equation}
The field $A_{\alpha0}$ is not an independent canonical variable, but rather
is defined in terms of the other variables by the constraint \eqref{eq:15.4.3}.
To see this, note that the `electric' field strengths $F_\alpha{}^{\mu0}$ are
\begin{equation}
  F_\alpha{}^{i0}=\Pi_{\alpha i},
  \qquad F_\alpha{}^{30}=\partial_3A_\alpha{}^0,
  \tag{15.4.6}\label{eq:15.4.6}
\end{equation}
so the constraint \eqref{eq:15.4.3} reads
\begin{equation}
  -(\partial_3)^2A_\alpha{}^0
  =\partial_i\Pi_{\alpha i}
  +\Pi_{\gamma i}C_{\gamma\alpha\beta}A_{\beta i}+J_\alpha{}^0,
  \tag{15.4.7}\label{eq:15.4.7}
\end{equation}
which can easily be solved (with reasonable boundary conditions) to give
$A_\alpha{}^0$ as a functional of $\Pi_{\gamma i}$, $A_{\beta i}$, and
$J_\alpha{}^0$. (We are using a summation convention, with indices $i$, $j$,
etc. summed over the values 1 and 2.) It should be noted that the canonical
conjugate to the matter field $\psi_\ell$ is
\begin{equation}
  \pi_\ell=\frac{\partial\mathcal L}{\partial(\partial_0\psi_\ell)}
  =\frac{\partial\mathcal L_M}{\partial(D_0\psi_\ell)},
  \tag{15.4.8}\label{eq:15.4.8}
\end{equation}
so the time component of the matter current can be expressed in terms of the
canonical variables of the matter fields \emph{alone}
\begin{equation}
  J_\alpha{}^0=-i\frac{\partial\mathcal L_M}{\partial D_0\psi_\ell}
  (t_\alpha)_{\ell m}\psi_m
  =-i\pi_\ell(t_\alpha)_{\ell m}\psi_m.
  \tag{15.4.9}\label{eq:15.4.9}
\end{equation}
Hence Eq.~\eqref{eq:15.4.7} defines $A_\alpha{}^0$ at a given time as a
functional of the canonical variables $\Pi_{\gamma i}$, $A_{\beta i}$,
$\pi_\ell$, and $\psi_m$ at the same time.

Now that we have identified the canonical variables in this gauge, we can
proceed to the construction of a Hamiltonian. The Hamiltonian density is
\begin{align}
  \mathcal H
  &=\Pi_{\alpha i}\partial_0A_{\alpha i}
  +\pi_\ell\partial_0\psi_\ell-\mathcal L
  \notag\\
 &=\Pi_{\alpha i}\bigl(F_{\alpha0i}+\partial_iA_{\alpha0}
 -C_{\alpha\beta\gamma}A_{\beta0}A_{\gamma i}\bigr)
  +\pi_\ell\partial_0\psi_\ell
  \notag\\
  &\quad-\frac12F_{\alpha0i}F_{\alpha0i}
  +\frac12F_{\alpha ij}F_{\alpha ij}
  +\frac12F_{\alpha i3}F_{\alpha i3}
  -\frac12F_{\alpha03}F_{\alpha03}-\mathcal L_M.
  \tag{15.4.10}\label{eq:15.4.10}
\end{align}
Using Eqs.~\eqref{eq:15.4.4} and \eqref{eq:15.4.6}, this is
\begin{align}
  \mathcal H={}&\mathcal H_M
 +\Pi_{\alpha i}\bigl(\partial_iA_{\alpha0}
 -C_{\alpha\beta\gamma}A_{\beta0}A_{\gamma i}\bigr)
  +\frac12\Pi_{\alpha i}\Pi_{\alpha i}
  \notag\\
  &+\frac12F_{\alpha ij}F_{\alpha ij}
  +\frac12\partial_3A_{\alpha i}\partial_3A_{\alpha i}
  -\frac12\partial_3A_{\alpha0}\partial_3A_{\alpha0},
  \tag{15.4.11}\label{eq:15.4.11}
\end{align}
where $\mathcal H_M$ is the matter Hamiltonian density:
\begin{equation}
  \mathcal H_M\equiv\pi_\ell\partial_0\psi_\ell-\mathcal L_M.
  \tag{15.4.12}\label{eq:15.4.12}
\end{equation}

Following the general rules derived in Section 9.2, we can now use this
Hamiltonian density to calculate matrix elements as path integrals over
$A_{\alpha i}$, $\Pi_{\alpha i}$, $\psi_\ell$, and $\pi_\ell$, with weighting
factor $\exp(iI)$, where
\begin{equation}
  I=\int d^4x\left[\Pi_{\alpha i}\partial_0A_{\alpha i}
  +\pi_\ell\partial_0\psi_\ell-\mathcal H+\epsilon\text{ terms}\right],
  \tag{15.4.13}\label{eq:15.4.13}
\end{equation}
in which the `$\epsilon$ terms' serve only to supply the correct imaginary
infinitesimal terms in propagator denominators. (See Section 9.2.) We note that
Eqs.~\eqref{eq:15.4.7} and \eqref{eq:15.4.9} give $A_\alpha{}^0$ as a
functional of the canonical variables, linear in $\Pi_{\alpha i}$ and
$\pi_\ell$. Inspection of Eq.~\eqref{eq:15.4.11} shows then (assuming
$\mathcal L_M$ to be no more than quadratic in $D_\mu\psi$) that the integrand
of the complete action \eqref{eq:15.4.13} is no more than quadratic in
$\Pi_{\alpha i}$ and $\pi_\ell$. We could therefore carry out the path
integral over these canonical `momenta' by the usual rules of Gaussian
integration. The trouble with this procedure is that the coefficients of the
terms in Eq.~\eqref{eq:15.4.13} of second order in $\Pi_{\alpha i}$ are
functions of the $A_{\alpha i}$, so the Gaussian integral would yield an
awkward field-dependent determinant factor. Also, the whole formalism at this
point looks hopelessly non-Lorentz-invariant.

Instead of proceeding in this way, we will apply a trick like that used in the
path integral formulation of electrodynamics in Section 9.6. Note that if for
a moment we think of $A_{\alpha0}$ as an independent variable, then the action
\eqref{eq:15.4.13} is evidently quadratic in $A_{\alpha0}$, with the
coefficient of the second-order term $A_{\alpha0}(x)A_{\beta0}(y)$ equal to the
field-independent kernel $(\partial_3)^2\delta^4(x-y)$. As we saw in the
appendix to Chapter 9, the integral of such a Gaussian over $A_{\alpha0}(x)$
is, up to a constant factor, equal to the value of the integrand at the
stationary `point' of the argument of the exponential. But the variational
derivative of the action here is
\begin{equation*}
  \frac{\delta I}{\delta A_{\alpha0}}
  =-\frac{\partial\mathcal H}{\partial A_{\alpha0}}
  =J_\alpha{}^0+\partial_i\Pi_{\alpha i}
  +C_{\beta\alpha\gamma}\Pi_{\beta i}A_{\gamma i}
  -\partial_3^2A_\alpha{}^0,
\end{equation*}
so the stationary `point' of the action is the solution of the constraint
equation \eqref{eq:15.4.7}. Hence, instead of using for $A_{\alpha0}$ the
solution of Eq.~\eqref{eq:15.4.7}, we can just as well treat it as an
independent variable of integration.

With $A_{\alpha0}$ now regarded as an independent variable, the Hamiltonian
$\int d^3x\,\mathcal H$ is evidently quadratic in $\Pi_{\alpha i}$, with the
coefficient of the second-order term $\Pi_{\alpha i}(x)\Pi_{\beta j}(y)$ given
by the field-independent kernel $\tfrac12\delta^4(x-y)\delta_{ij}$. Assuming
that the same is true for the matter variable $\pi_\ell$, we can evaluate path
integrals over $\pi_\ell$ and $\Pi_{\alpha i}$ up to a constant factor by
simply setting $\pi_\ell$ and $\Pi_{\alpha i}$ at the stationary `points' of
the action corresponding to Eq.~\eqref{eq:15.4.1}:
\begin{align*}
  0&=\frac{\delta I}{\delta\pi_\ell}
  =\partial_0\psi_\ell-\frac{\partial\mathcal H_M}{\partial\pi_\ell},\\
  0&=\frac{\delta I}{\delta\Pi_{\alpha i}}
  =\partial_0A_{\alpha i}-\Pi_{\alpha i}-\partial_iA_{\alpha0}
  +C_{\alpha\beta\gamma}A_{\beta0}A_{\gamma i}
  =F_{\alpha0i}-\Pi_{\alpha i}.
\end{align*}
Inserting these back into Eq.~\eqref{eq:15.4.13} gives
\begin{align}
  I&=\int d^4x\left[\mathcal L_M+\frac12F_{\alpha0i}F_{\alpha0i}
  -\frac12F_{\alpha ij}F_{\alpha ij}
  -\frac12\partial_3A_{\alpha i}\partial_3A_{\alpha i}
  +\frac12(\partial_3A_{\alpha0})^2\right]
  \notag\\
  &=\int d^4x\,\mathcal L,
  \tag{15.4.14}\label{eq:15.4.14}
\end{align}
where $\mathcal L$ is the Lagrangian \eqref{eq:15.3.1} with which we started!
In other words, \emph{we are to do path integrals over $\psi_\ell(x)$ and all
four components of $A_{\alpha\mu}(x)$, with a manifestly covariant weighting
factor $\exp(iI)$ given by Eqs.~\eqref{eq:15.4.14} and \eqref{eq:15.3.1}, but
with the axial-gauge condition} enforced by inserting a factor
\begin{equation}
  \prod_{x,\alpha}\delta\bigl(A_{\alpha3}(x)\bigr).
  \tag{15.4.15}\label{eq:15.4.15}
\end{equation}
As long as $\mathcal O_A$, $\mathcal O_B$, $\ldots$ are gauge-invariant, we have
\begin{align}
  \bigl\langle T\{\mathcal O_A\mathcal O_B\cdots\}\bigr\rangle_{\mathrm{VACUUM}}
  \propto{}&\int\left[\prod_{\ell,x}d\psi_\ell(x)\right]
  \left[\prod_{\alpha,\mu,x}dA_{\alpha\mu}(x)\right]
  \notag\\
  &\cdot\mathcal O_A\mathcal O_B\cdots
  \exp\{iI+\epsilon\text{ terms}\}
  \prod_{x,\alpha}\delta\bigl(A_{\alpha3}(x)\bigr),
  \tag{15.4.16}\label{eq:15.4.16}
\end{align}
with Lorentz- and gauge-invariant action $I$ given by Eq.~\eqref{eq:15.4.14}.

\begin{center}
  * * *
\end{center}

For future reference, we note that the volume element
$\prod_{\alpha,\mu,x}dA_{\alpha\mu}(x)$ for the integration over gauge fields
in \eqref{eq:15.4.16} is gauge-invariant, in the sense that
\begin{equation}
  \prod_{\alpha,\mu,x}dA_{\Lambda\,\alpha\mu}(x)
  =\prod_{\alpha,\mu,x}dA_{\alpha\mu}(x),
  \tag{15.4.17}\label{eq:15.4.17}
\end{equation}
where $A_{\Lambda\,\alpha\mu}(x)$ is the result of acting on
$A_{\alpha\mu}(x)$ with a gauge transformation having transformation
parameters $\Lambda_\alpha(x)$. It will be enough to show that this is true for
transformations near the identity, say with infinitesimal transformation
parameters $\lambda_\alpha(x)$. In this case,
\begin{equation*}
  A_{\lambda\,\alpha}{}^\mu=A_\alpha{}^\mu+\partial^\mu\lambda_\alpha
  +C_{\alpha\beta\gamma}A_\beta{}^\mu\lambda_\gamma,
\end{equation*}
so the volume elements are related by
\begin{equation*}
  \prod_{\alpha,\mu,x}dA_{\lambda\,\alpha\mu}(x)
  =\operatorname{Det}(\mathcal N)
  \prod_{\alpha,\mu,x}dA_{\alpha\mu}(x),
\end{equation*}
where $\mathcal N$ is the `matrix'
\begin{equation*}
  \mathcal N_{\alpha\mu x,\beta\nu y}
  \equiv\frac{\delta A_{\lambda\,\alpha\mu}(x)}
  {\delta A_{\beta\nu}(y)}
  =\delta^4(x-y)\delta_\mu{}^\nu
  \bigl[\delta_{\alpha\beta}+C_{\alpha\beta\gamma}\lambda_\gamma(x)\bigr].
\end{equation*}
The determinant of $\mathcal N$ is unity to first order in $\lambda_\gamma$,
because the trace $C_{\alpha\alpha\gamma}$ vanishes.

In this chapter we shall assume that the volume element
$\prod_{n,x}d\psi_n(x)$ for the integration over matter fields is also
gauge-invariant. There are important subtleties here, to which we shall return
in Chapter 22, but as shown there this assumption turns out to be valid in our
present non-Abelian gauge theories of strong and electroweak interactions.
